\documentclass[border=6pt]{standalone}
\usepackage{amsmath}
\usepackage{amssymb}
\usepackage{array}
\usepackage{booktabs}
\begin{document}
\begin{tabular}{@{} >{\bfseries\raggedright\arraybackslash}p{1.6cm} >{\centering\arraybackslash}p{4.6cm} >{\centering\arraybackslash}p{4.6cm} >{\centering\arraybackslash}p{4.6cm} @{}}
\toprule
\multicolumn{4}{c}{\textbf{Function Operations}} \\
\midrule
 & \textbf{Add or Subtract Functions} & \textbf{Multiply or Divide Functions} & \textbf{Compose Functions} \\
\midrule
ALGEBRA &
\(\begin{gathered}(f+g)(x)=f(x)+g(x)\\ (f-g)(x)=f(x)-g(x)\end{gathered}\) &
\(\begin{gathered}(f\cdot g)(x)=f(x)\cdot g(x)\\ \left(\dfrac{f}{g}\right)(x)=\dfrac{f(x)}{g(x)}\end{gathered}\) &
\(\begin{gathered}(f\circ g)(x)=f(g(x))\\ (g\circ f)(x)=g(f(x))\end{gathered}\) \\
\midrule
WORDS &
The domain of the sum or difference of $f$ and $g$ is the intersection of the domain of $f$ and the domain of $g$. &
The domain is the set of all real numbers for which $f$ and $g$ and the new function are defined. &
The domain of $f\circ g$ is the set of all real numbers $x$, in the domain of $g$, such that $g(x)$ is in the domain of $f$. \\
\midrule
NUMBERS &
\(\begin{gathered}\text{Let } f(x)=3x+5 \\ \text{and } g(x)=x-3 \\ f+g=(3x+5)+(x-3)=4x+2 \\ f-g=(3x+5)-(x-3)=2x+8\end{gathered}\) &
\(\begin{gathered}\text{Let } f(x)=3x+5 \\ \text{and } g(x)=x-3 \\ f\cdot g=(3x+5)(x-3)=3x^{2}-4x-15 \\ \dfrac{f}{g}=\dfrac{3x+5}{x-3} \text{ for } x\ne 3\end{gathered}\) &
\(\begin{gathered}\text{Let } f(x)=3x+5 \\ \text{and } g(x)=x-3 \\ f\circ g=3(x-3)+5=3x-4 \\ g\circ f=(3x+5)-3=3x+2\end{gathered}\) \\
\bottomrule
\end{tabular}
\end{document}
